\documentclass[12pt]{article}
\usepackage[margin=1in]{geometry}
\usepackage{amsmath}
\usepackage{amssymb}
\usepackage{array}
\usepackage{booktabs}
\usepackage{graphicx}

\title{\textbf{GED Math Worksheet: Tables, Graphs, and Interpreting Data}}
\author{}
\date{}

\begin{document}

\maketitle

\noindent\textbf{Name:} \underline{\hspace{3in}} \quad \textbf{Date:} \underline{\hspace{2in}}

\vspace{0.2in}

\noindent\textbf{Instructions:} Answer all 30 questions. Show your work where applicable. You may use a calculator for these problems.

\vspace{0.3in}

\section*{Part A: Tables (Questions 1-6)}

\noindent\textbf{Questions 1-3 refer to the following table:}

\begin{center}
\begin{tabular}{|c|c|c|c|c|}
\hline
\textbf{Month} & \textbf{Jan} & \textbf{Feb} & \textbf{Mar} & \textbf{Apr} \\
\hline
\textbf{Sales (\$1000s)} & 45 & 52 & 38 & 61 \\
\hline
\end{tabular}
\end{center}

\begin{enumerate}
\item What were the total sales for the four months?
\vspace{0.8in}

\item What is the average monthly sales for this period?
\vspace{0.8in}

\item By how much did sales increase from February to April?
\vspace{0.8in}
\end{enumerate}

\noindent\textbf{Questions 4-6 refer to the following table:}

\begin{center}
\begin{tabular}{|c|c|c|c|}
\hline
\textbf{Product} & \textbf{Price} & \textbf{Quantity Sold} & \textbf{Revenue} \\
\hline
\textbf{Widget A} & \$8 & 125 & \$1,000 \\
\hline
\textbf{Widget B} & \$12 & 87 & \$1,044 \\
\hline
\textbf{Widget C} & \$15 & 92 & \$1,380 \\
\hline
\textbf{Widget D} & \$10 & 156 & \$1,560 \\
\hline
\end{tabular}
\end{center}

\begin{enumerate}
\setcounter{enumi}{3}
\item Which product generated the highest revenue?
\vspace{0.8in}

\item If Widget B's price increases by 20\%, what is the new price?
\vspace{0.8in}

\item What is the total number of widgets sold across all products?
\vspace{0.8in}
\end{enumerate}

\section*{Part B: Bar Graphs (Questions 7-12)}

\noindent\textbf{Questions 7-9 refer to the following data about monthly temperatures (in °F):}

\noindent Month: Jan (32), Feb (35), Mar (45), Apr (58), May (72), Jun (85)

\begin{enumerate}
\setcounter{enumi}{6}
\item What is the difference in temperature between the warmest and coldest month listed?
\vspace{0.8in}

\item What is the average temperature for the six months shown?
\vspace{0.8in}

\item Between which two consecutive months is the greatest temperature increase?
\vspace{0.8in}
\end{enumerate}

\noindent\textbf{Questions 10-12 refer to the following data about student enrollment in classes:}

\noindent Math (120), Science (95), English (110), History (85), Art (60)

\begin{enumerate}
\setcounter{enumi}{9}
\item How many more students are enrolled in Math than in Art?
\vspace{0.8in}

\item What is the total enrollment across all five classes?
\vspace{0.8in}

\item What percentage of total enrollment is in Science? (Round to the nearest whole percent)
\vspace{0.8in}
\end{enumerate}

\section*{Part C: Line Graphs (Questions 13-17)}

\noindent\textbf{Questions 13-17 refer to the following data showing website traffic (in thousands of visitors) over six weeks:}

\noindent Week 1 (12), Week 2 (15), Week 3 (14), Week 4 (18), Week 5 (22), Week 6 (20)

\begin{enumerate}
\setcounter{enumi}{12}
\item During which week did the website experience the highest traffic?
\vspace{0.8in}

\item What was the change in visitors from Week 1 to Week 6?
\vspace{0.8in}

\item During which consecutive weeks was there a decrease in visitors?
\vspace{0.8in}

\item What is the average number of visitors across all six weeks?
\vspace{0.8in}

\item If the trend from Week 4 to Week 5 continues, how many visitors would you predict for Week 7?
\vspace{0.8in}
\end{enumerate}

\section*{Part D: Circle Graphs (Questions 18-22)}

\noindent\textbf{Questions 18-22 refer to the following data about a budget breakdown for a household (3,600 dollars total):}

\noindent Housing (40\%), Food (25\%), Transportation (18\%), Utilities (12\%), Entertainment (5\%)

\begin{enumerate}
\setcounter{enumi}{17}
\item How much money is allocated to Housing?
\vspace{0.8in}

\item What is the difference between the amount spent on Food and Entertainment?
\vspace{0.8in}

\item If the Entertainment budget is increased by 3\%, how much money would be allocated to it?
\vspace{0.8in}

\item What percentage of the budget is spent on Housing and Transportation combined?
\vspace{0.8in}

\item If total household income increases to \$4,000 but percentages remain the same, how much will be spent on Food?
\vspace{0.8in}
\end{enumerate}

\section*{Part E: Scatter Plots (Questions 23-26)}

\noindent\textbf{Questions 23-26 refer to the following data showing the relationship between study hours and test scores:}

\noindent Study Hours: 1, 2, 3, 4, 5, 6
\noindent Test Scores: 62, 68, 75, 81, 87, 92

\begin{enumerate}
\setcounter{enumi}{22}
\item What type of correlation does this data show?
\vspace{0.8in}

\item Based on the pattern, what would you estimate the test score to be for 7 hours of study?
\vspace{0.8in}

\item What is the range of test scores shown in the data?
\vspace{0.8in}

\item Is there an outlier in this dataset? Explain your answer.
\vspace{0.8in}
\end{enumerate}

\section*{Part F: Histograms (Questions 27-29)}

\noindent\textbf{Questions 27-29 refer to the following histogram data showing distribution of ages at a community center:}

\noindent Age 18-25 (8), Age 26-35 (15), Age 36-45 (12), Age 46-55 (10), Age 56-65 (5)

\begin{enumerate}
\setcounter{enumi}{26}
\item What is the total number of people represented in the histogram?
\vspace{0.8in}

\item Which age group has the highest frequency?
\vspace{0.8in}

\item What percentage of people are age 46 or older? (Round to the nearest whole percent)
\vspace{0.8in}
\end{enumerate}

\section*{Part G: Box Plots (Question 30)}

\noindent\textbf{Question 30 refers to the following dataset: 12, 15, 18, 22, 25, 28, 30, 32, 35}

\begin{enumerate}
\setcounter{enumi}{29}
\item Find the median, lower quartile, upper quartile, and interquartile range for this dataset.
\vspace{1.5in}
\end{enumerate}

\vspace{1in}

\noindent\textbf{End of Worksheet}

\end{document}